\documentclass[10pt,a4paper]{amsart}
\usepackage[margin=19mm]{geometry}
\usepackage{amsmath,amssymb,mathtools}
\usepackage[scr=rsfs,cal=euler]{mathalfa}
\usepackage{xfrac}
\usepackage[unicode,hidelinks,bookmarks=false]{hyperref}
\newcommand{\ie}{i.e.\ }
\newcommand{\cf}{cf.\ }
\newcommand{\resp}{respectively\ }
\newcommand{\Subsection}{\subsection}
\providecommand{\nobreakdash}{\nobreak\hbox{-}}
\setlength{\emergencystretch}{2em}
\multlinegap0pt
\usepackage{mathtools}
\usepackage{amssymb}
\newtheorem{lemma}{Lemma}
\title{\textbf{Quasi-recursive MDS Matrices over Galois Rings}}
\author{Shakir Ali, Atif Ahmad Khan, Abhishek Kesarwani, Susanta Samanta}
\date{Source: arXiv:2512.17256v1 (CC BY 4.0)}
\begin{document}
\maketitle

\noindent\emph{Source and licence.} \textbf{Quasi-recursive MDS Matrices over Galois Rings}. Version arXiv:2512.17256v1 (CC BY 4.0), distributed by arXiv under the Creative Commons Attribution 4.0 International licence, http://creativecommons.org/licenses/by/4.0/, as recorded in the arXiv licence metadata for this exact version and shown on the abstract page https://arxiv.org/abs/2512.17256v1. Full licence notice: \texttt{licenses/arxiv-2512.17256-CC-BY-4.0.txt}.\par\smallskip

\noindent\textit{Editorial scope.} Counted unit: the article's lemma companiontocyclic (source0.tex lines 549--561). The article's complete original proof of it is delivered with it (source0.tex lines 562--656, one \texttt{proof} environment of 95 lines). No mathematics is written, repaired or re-derived: the statement and the proof are the article's own lines, in the article's own notation, and no retained line is substituted or re-flowed.  No further source-local context is needed: every cross-reference of every retained line resolves inside the two delivered spans.  Nothing was omitted from any retained span, so the package carries no omission record.  No retained line contains a citation, so the wrapper carries no bibliography and no bibliography entry.  No retained line contains a figure environment, an image-inclusion command or drawing code, so no image is delivered and the package declares no asset dependency.  Editorial formatting only, in addition: an AMS \texttt{amsart} preamble that restates the article's own declarations and macros used by the retained lines, namely the environments \texttt{lemma} on separate counters, together with the standard packages those lines need.  The retained environments print in this document as Lemma 1 in order of appearance, whereas the article prints the same items with the numbering of its own full document.  The complete native source of the article as distributed in the arXiv e-print for this exact version is bundled unchanged as \texttt{sources/191303/source0.tex}.
\begin{lemma}\label{companiontocyclic}
    Let $g(X)=g_0+g_1X+\cdots+g_{k-1}X^{k-1}+X^k\in GR(p^s,p^{sm})[X;\sigma]$ be a right divisor of $X^n-1.$ 
    Then, for any positive integer $t$,
    \[
    C^{[t-1]}_g \cdot C^{[t-2]}_g \cdots C^{[1]}_g \cdot C_g =
    \begin{bmatrix}
        X^t \bmod_\ast g\\
        X^{t+1} \bmod_\ast g\\
        \vdots\\
        X^{t+k-1} \bmod_\ast g
    \end{bmatrix}_{k\times k}.
    \]
\end{lemma}

\begin{proof}
    	We show by strong induction that for all integers $i\geq 1$, the following property $P_i$ is satisfied:
	\begin{eqnarray}\label{eqn201}
		P_i&:&\begin{bmatrix}
			X^i ~ \mod _*g \\
			X^{i+1} ~ \mod _*g  \\
			\vdots\\
			X^{i+k-1} ~ \mod _*g  \\
		\end{bmatrix}_{k\times k}=C^{[i-1]}_g\cdot C^{[i-2]}_g\cdots C^{[1]}_g\cdot C_g.
	\end{eqnarray}	
    
    For $i=1,$ we have 
	\begin{eqnarray}
		P_0&:&\begin{bmatrix}
			X^1 ~ \mod _*g  \\
			X^{2} ~ \mod _*g  \\
			\vdots\\
			X^{k} ~ \mod _*g  \\
			
		\end{bmatrix}_{k\times k}=C_g.
	\end{eqnarray}
    
	Thus, (\ref{eqn201}) is true for $i=1$. Suppose our assumption is true for $1,~2,~\dots,~i-1$. Then, we have
	
	%%%%%%%%%%%%%%%%%%%%%%%%%%%%%%%%%%%%%%%%%%%%%%%$\binom nk=^nC_k=\frac{n!}{k!(n-k)!}$
	
	$$\begin{bmatrix}\label{eqn5.8}
		X^i ~ \mod _*g  \\
		X^{i+1} ~ \mod _*g  \\
		\vdots\\
		X^{i+k-1} ~ \mod _*g  \\
	\end{bmatrix}_{k\times k}=
	C_g^{[i-1]}\cdot C_g^{[i-2]}\cdots C_g^{[1]}\cdot C_g.$$
	Now, we prove (\ref{eqn201}) for $i$. For this, we compute 
    \begin{eqnarray*}
        \notag&~&C^{[i]}_g C_g^{[i-1]}\cdot C_g^{[i-2]}\cdots C_g^{[1]}\cdot C_g\\
		\notag&~~~~~&	=\begin{bmatrix}
			0&1&0&\cdots&0\\
			0&0&1&\cdots&0\\
			\vdots&\vdots&\vdots&\ddots&\vdots&\\
			0&0&0&\cdots&1\\
			\sigma^{i}(g_0)&\sigma^{i}(g_1)&\sigma^{i}(g_2)&\cdots&\sigma^{i}(g_{k-1})
		\end{bmatrix}_{k\times k} \cdot \begin{bmatrix}
			X^i ~ \mod _*g  \\
			X^{i+1} ~ \mod _*g \\
			\vdots\\
			X^{i+k-1} ~ \mod _*g \\
		\end{bmatrix}_{k\times k}\\
    \end{eqnarray*}
	\begin{eqnarray}\label{eqn11}
		\notag&=&\begin{bmatrix}
			X^{i+1}~ \mod _*g  \\
			X^{i+2}~ \mod _*g  \\
			\vdots\\
			X^{i+k-1}~\mod_*g\\
            \sigma^{i}(g_0)X^i+\cdots+\sigma^{i}(g_{k-1})X^{i+k-1}~ \mod _*g  
		\end{bmatrix}_{k\times k}\\
		&=&\begin{bmatrix}
			X^{i+1}~ \mod _*g \\
			X^{i+2}~ \mod _*g  \\
			\vdots\\
			\Big(\sum_{j=0}^{k-1}\sigma^{i}(g_j)X^{j+i}\Big)~ \mod _*g  \\
		\end{bmatrix}_{k\times k}.
	\end{eqnarray}
	\normalsize
    
	Since,
	\begin{eqnarray*}
		X^i*g_t&=&\sigma^i(g_t)X^i,
	\end{eqnarray*}
    
	we have
	\begin{eqnarray}\label{eqn12}
		X^{i+k}=X^i\ast X^k\notag&=& X^i*(g_0+g_1X+g_2X^2+\cdots+g_{k-1}X^{k-1})\\
\notag&=&\sigma^i(g_0)X^i+\sigma^i(g_1)X^{i+1}+\cdots+\sigma^i(g_{k-1})X^{i+k-1}\\
		&=&\sum_{j=0}^{k-1}\sigma^i (g_j)X^{i+j}.
	\end{eqnarray}
    
	From the last row of the matrix in (\ref{eqn11}), and (\ref{eqn12}), we obtain
	\begin{eqnarray}\label{eqn13}
		\Big(\sum_{j=0}^{k-1}\sigma^{i}(g_j)X^{j+i}\Big)~ \mod _*g  &=&X^{i+k}~ \mod _*g.
	\end{eqnarray}

Combining (\ref{eqn11}) and  (\ref{eqn13}), it follows that
	\begin{eqnarray}\label{eqn202}
		C_g^{[i]}\cdot 	C_g^{[i-1]}\cdot C_g^{[i-2]}\cdots C_g^{[1]}\cdot C_g=\begin{bmatrix}
			X^{i+1} ~ \mod _*g  \\
			X^{i+2} ~ \mod _*g  \\
			\vdots\\
			X^{i+k} ~ \mod _*g  \\
		\end{bmatrix}_{k\times k}.
	\end{eqnarray}
    
	Hence, $P_i$ is true for all integer $i\geq 1.$
\end{proof}

\end{document}
